% 07-conclusoes.tex -- incluído por main-pt.tex
% =====================================================================
\section{Conclusões}
\label{sec:conclusion}
% =====================================================================

Esta contribuição investigou a viabilidade do consenso bizantino em um registro nacional de créditos de carbono baseado em blockchain. A questão foi formulada em termos de duas escalas distintas: a escala de governança, que determina o número de validadores e o custo de comunicação de cada rodada de consenso, e a escala de mercado, que determina a demanda de transações e sua concentração em torno dos prazos de conciliação. Os protocolos avaliados foram selecionados a partir de requisitos do domínio, o tráfego de consenso foi medido em nível de pacote, validadores bizantinos foram incluídos no projeto fatorial, e as elasticidades em relação às duas escalas foram estimadas por metamodelos validados fora da faixa de ajuste, com as interações testadas por ANOVA. Mostrou-se que \tbd{principal achado sobre as elasticidades e a restrição vinculante} e que \tbd{fronteiras de viabilidade na faixa de governança, com e sem falhas}, o que \tbd{confirma/refuta} a hipótese de que a viabilidade de um registro nacional depende mais do desenho de governança do que do volume de transações. Em qualquer caso, as fronteiras obtidas podem subsidiar a regulamentação do registro central do SBCE, ao explicitar quantos validadores um registro pode comportar para cada nível de demanda. Possíveis extensões deste trabalho envolvem a simulação de arquiteturas hierárquicas de consenso, a inclusão de outros comportamentos bizantinos e a calibração do modelo de demanda com os dados dos primeiros períodos de compromisso do SBCE. Outra direção de pesquisa poderia investigar o impacto dos oráculos de MRV sobre a demanda e a segurança do registro.
